%% GENERATED by python/paper/build.py from results/ -- do not edit by hand.
%% Rebuild:  .venv\Scripts\python.exe python\paper\build.py --only US_Ageing_vectorX
%% Source:   results/shocks/US_shocks.csv
\begin{table}[!htb]
\centering
\begin{threeparttable}
\caption{The effect of ageing in US -- 2020 (vector $X_i$ calibration)}
\label{table:US:ageing_vectorX}
\renewcommand{\arraystretch}{1.25}
\begin{tabularx}{.9\textwidth}{p{3cm}YYY}
\toprule
\textbf{Scenario} & \textbf{Tax rate} & \textbf{Savings rate} & \textbf{Avg. workweek} \\
\midrule 
Baseline & 14.43\% & 15.37\% & 39.39 \\[1.25ex]
\multicolumn{4}{l}{\textit{Full effect:}} \\\hline
Mild ageing\tnote{a} & 18.45\% & $-1.00$ p.p. & 39.08 \\
Acute ageing\tnote{b} & 23.11\% & $-2.08$ p.p. & 38.66 \\[1.25ex]
\multicolumn{4}{l}{\textit{Economic equilibrium effect:}} \\\hline
Mild ageing\tnote{a} & 14.43\% & 0.00 p.p. & 40.12 \\
Acute ageing\tnote{b} & 14.43\% & 0.00 p.p. & 40.99 \\[1.25ex]
\bottomrule
\end{tabularx}
\begin{tablenotes}
\footnotesize
\item Each scenario is a separate equilibrium path: the demography holds throughout and the economy starts from its own steady state, so the capital stock brought into 2020 is the counterfactual one rather than the baseline\textquotesingle s. The savings rate is savings relative to GDP; the baseline row reports its level and every other row the change against that baseline in percentage points.
\item[a] The ``mild ageing'' scenario refers to the case with $\nu_t$ set at $(1+\nu_t^{base})/2$ throughout.
\item[b] The ``acute ageing'' scenario refers to $\nu_t = 1$ throughout. Vector-$X_i$ calibration: $X_i$ is identified from relative hours, which are data here, and the level of $ar h$ is then not identified --- only its ratio to the baseline is. $eta$, $\omega$, the tax rate and the savings rate are common to the two variants; what differs is anything defined through $\eta$ or $X$.
\end{tablenotes}
\end{threeparttable}
\end{table}
